% ==================================================================
% experiments.tex -- computational study on five UA-FLP benchmark instances
%
% Every number is a macro from numbers.tex (written by make_tables.py from the result files).
% Paths: \bxdir is the folder that holds tables/ and figures/ (empty when compiled from paper/).
% In the main paper set e.g. \newcommand{\bxdir}{benchmark_v1/} before \input-ing this file.
% ==================================================================
\providecommand{\bxdir}{}
\input{\bxdir numbers}

\section{Computational experiments}
\label{sec:bx}

This section compares the two-stage method with five heuristic baselines and an exact convex MINLP on five
unequal-area facility layout (UA-FLP) instances from the literature. All methods solve the same model:
rectangular departments with exact areas, a maximum aspect ratio $\alpha$, a fixed rectangular floor, no
overlap, and the objective $\sum_{i<j} w_{ij}\,(|x_i-x_j|+|y_i-y_j|)$ on department centroids.

\subsection{Test instances}
\label{sec:bx-instances}

\Cref{tab:bx-setup} lists the instances. O7, vC10a, AB20 and Du62 fill the floor completely (Du62 up to
rounding of the floor side), so every feasible layout is a partition of the floor. SC30 leaves
\bxFreeSCthirty\,\% of the floor free; empty space is allowed in our model.
The floors and aspect limits are those used by Jeong and Chae~\cite{jeong2026}; for Du62 we use the
$117.124\times117.124$ square on which they report their results. O7 is also instance \texttt{o7} of
MINLPLib~\cite{minlplib_o7}, which takes it from Westerlund and Lundqvist~\cite{westerlund2001}; MINLPLib lists
the global optimum \bxRefOsevenFull, with dual bounds of SCIP and BARON equal to this value. Two data details
matter for the comparison with published values. First, vC10a uses $w_{9,10}=59.2$; this is the value under which
the layout reported in~\cite{jeong2026} has cost 18{,}522.79. Second, the AB20 flow matrix is the one of the Kang
and Chae data file, which equals the matrix of Armour and Buffa~\cite{armour1963} multiplied by 10 (with the area
and floor scaled accordingly).

\input{\bxdir tables/tab_setup}

\subsection{Methods and settings}
\label{sec:bx-methods}

All stochastic baselines, and Stage~1 of our method, work on the slicing-tree (STS) encoding
of Kang and Chae~\cite{kang2017} and Jeong and Chae~\cite{jeong2026}: a department sequence, a cut order and
a cut code per cut. Cuts split the current rectangle in proportion to the areas of the two parts, so areas are
exact. For SC30 the free area is split into 17 dummy blocks of equal area, without flow and without a shape
limit, as in the Kang and Chae data file. The STS methods use the same static penalty
$\pi=0.05\,(\sum_{i<j}w_{ij})(W+H)$ per unit of aspect excess, and they return the best layout that meets
the aspect limit exactly. Slicing layouts cannot represent every layout; on O7 the optimal layout is not a slicing
layout.

\begin{description}[leftmargin=0pt]
\item[Construct.] A deterministic construction heuristic. Departments are ordered greedily by flow to the
departments already ordered (a CORELAP-style rule on the flow matrix), then placed in a flexible bay
structure. All bay counts $k=1,\dots,n$ and both bay orientations are tried; the best layout that meets the
aspect limit is returned. One run per instance.
\item[Swap.] CRAFT-style pairwise exchange~\cite{armour1963} on the STS sequence: steepest descent over all
pairwise swaps, restarted from a random genome at each local optimum.
\item[GA.] A genetic algorithm without local search: population 50, two elites, binary tournament, PMX on
the sequence and the cut order, one-point crossover on the cut codes, mutation probability 0.2.
\item[SA.] Simulated annealing with three move types; the start temperature accepts a median worsening move
with probability 0.75, and the temperature falls geometrically in elapsed time to $10^{-4}$ of its start value.
\item[CRO.] Our reproduction of the coral reefs optimisation of Jeong and Chae~\cite{jeong2026} (grid
$7\times7$, parameters of their Table~5, duplicate larvae rejected). The paper's stopping rule is replaced by the
common time budget. Run to the paper's stopping rule, our reproduction reaches the published best value on O7 in
\bxCroFullOsevenHits{} runs and on vC10a in \bxCroFullvCtenaHits{} runs; on AB20 ($7\times7$) its mean is
\bxCroFullOurMeanABtwenty{} against the published \bxCroFullPaperMeanABtwenty, on Du62 its mean is
\bxCroFullOurMeanDusixtytwo{} against \bxCroFullPaperMeanDusixtytwo, and on SC30 it is not reproduced
(mean \bxCroFullOurMeanSCthirty{} against \bxCroFullPaperMeanSCthirty). ``CRO'' below always means our
reproduction.
\item[SA + polish, CRO + polish.] A fairness control: the final layout of each SA and CRO run is converted
into a sequence pair and solved once by the Stage-2 SOCP of our method (the same model and margins). The SOCP
cannot make a layout worse; if it gives no gain beyond $10^{-8}$ relative, the original layout is kept. The polish
took at most \bxPolishTimeMaxSAPolish~s (SA) and \bxPolishTimeMaxCROPolish~s (CRO) per run and is not counted in
the budget. These rows separate the gain from the SOCP geometry from the gain from a better search.
\item[MINLP.] An exact convex mixed-integer nonlinear model solved by SCIP~10 (one thread, 300~s): continuous
positions and sizes, the area constraint $h_i\ge A_i/w_i$ (convex), linear aspect constraints, four binaries per
department pair with big-$M$ disjunctions, and absolute-value distances. Its optimum is the UA-FLP optimum (up to
the solver feasibility tolerance $10^{-7}$), and its dual bound is a valid lower bound.
\item[Ours.] The two-stage method of \cref{sec:method}. Stage~1 is slicing-tree annealing for half of the
budget and keeps an elite pool of up to four layouts. Stage~2 converts each elite layout into a sequence
pair~\cite{murata1996,liu2007} and searches over sequence pairs; each sequence pair is evaluated by a convex
second-order cone program (SOCP) solved by Clarabel~\cite{goulart2024}. If fewer than 3000 SOCP solves fit in the
remaining time, Stage~1 continues to the end of the budget and Stage~2 is a single SOCP solve (polish).
\end{description}

We also ran two restricted exact models, one run of 300~s each: a MIP in which every department chooses one of at
most nine fixed exact-area shapes (SCIP through OR-Tools~\cite{ortools}), and a CP-SAT model~\cite{cpsat} on an
integer grid. Both are infeasible by construction on the full floors (no tiling with so few shapes; grid rounding
leaves fewer cells than the areas need on O7 and Du62), and neither found a feasible layout on any instance. We do
not report them further.

The parameters of Swap, GA, SA and CRO are our own choices or the published ones and were not tuned on these
instances. The Stage-1 cycle length and temperatures of our method were tuned on AB20 and SC30, and the Stage-2
temperatures on SC30; both are test instances (\cref{sec:bx-threats}).

\subsection{Protocol}
\label{sec:bx-protocol}

\paragraph{Budgets and seeds.} Every stochastic method gets the same wall-clock budget per run
(\cref{tab:bx-setup}): \bxBudgetOseven~s for O7 and vC10a, \bxBudgetABtwenty~s for AB20 and SC30, and
\bxBudgetDusixtytwo~s for Du62. Each stochastic method is run with seeds 0--9. The construction heuristic and
the exact models are run once.

\paragraph{Machines.} All stochastic methods, including all runs of our method and its variants, ran on the same
machine: a Linux (aarch64) virtual machine with 4 cores, Python 3.10.12, NumPy 2.2.6 and Numba 0.68.0, with three
runs in parallel, each as a separate process. The exact models (MINLP, MIP, CP-SAT) ran on a different machine with 2 cores
(Intel Xeon, 2.8~GHz), one run at a time. The polish of the SA and CRO layouts was applied afterwards to the stored
layouts, and its time is recorded separately.

\paragraph{Budget overruns.} The runs of our method and its variants were made before a fix of the time check
in Stage~1 and can exceed the budget. The largest overrun is \bxOverOursMaxSOseven~s (\bxOverOursMaxPctOseven\,\%) on
O7, \bxOverOursMaxSvCtena~s (\bxOverOursMaxPctvCtena\,\%) on vC10a, \bxOverOursMaxSABtwenty~s
(\bxOverOursMaxPctABtwenty\,\%) on AB20, \bxOverOursMaxSSCthirty~s (\bxOverOursMaxPctSCthirty\,\%) on SC30 and
\bxOverOursMaxSDusixtytwo~s (\bxOverOursMaxPctDusixtytwo\,\%) on Du62. The baselines exceed the budget by at most
\bxOverBaseMaxS~s.

\paragraph{Independent checker and tolerance.} No method reports its own cost. After each run, the layout is
checked by an independent evaluator, which recomputes the cost from the rectangles and checks every constraint
with tolerance $\varepsilon=10^{-6}$: $w_ih_i\ge A_i(1-\varepsilon)$, aspect ratio at most $\alpha+\varepsilon$,
containment in the floor up to $\varepsilon$, and no pairwise overlap larger than $\varepsilon$ in both
directions. A run counts as feasible only if the evaluator accepts its layout. A second, separate verification
script re-checked all stored layouts: layouts produced by a solver (SOCP, MINLP) meet the limits to about $10^{-7}$,
and slicing layouts are exact to about $10^{-9}$. All \bxNlayouts{} stored layouts of the protocol were checked
again at the end; stored and recomputed cost and feasibility agree for every file. The code of each round of
runs is identified by MD5 checksums, and the test suite of the evaluator and models passed
(\bxTestsPassed{} tests passed, \bxTestsXfail{} expected failures).

\paragraph{Smoke tests.} Before the main runs, every heuristic was run for a few seconds on O7, vC10a and AB20, and
all produced feasible layouts. The construction heuristic first had no feasible layout on AB20 with vertical bays
only, which is why horizontal bays were added. The MINLP model was validated on O7 with a 3600~s run: SCIP proved
optimality at \bxMinlpOsevenProofT~s with value \bxMinlpOsevenProofVal, which matches the MINLPLib optimum up to the
solver tolerance.

\paragraph{Reference values and statistics.} Gaps are $100\,(\text{cost}-\text{ref})/\text{ref}$ with the
reference values in \cref{tab:bx-main}: O7 \bxRefOsevenFull{} (proven optimum, MINLPLib~\cite{minlplib_o7});
vC10a \bxRefvCtena{} and SC30 \bxRefSCthirty{} from Jeong and Chae~\cite[Table~6]{jeong2026}; AB20 \bxRefABtwenty{}
from Gon\c{c}alves and Resende~\cite[Table~12]{goncalves2015}; Du62 \bxRefDusixtytwo{} from Kang and
Chae~\cite{kang2017} as listed in~\cite[Table~6]{jeong2026}. \Cref{sec:bx-records} discusses how far each value is
comparable. To compare methods we use two-sided Mann--Whitney U tests on the ten costs per method (the seeds are
independent random streams, so the samples are unpaired), with Holm's correction over the comparisons within an
instance.

\subsection{Results}
\label{sec:bx-results}

\Cref{tab:bx-main} gives the results, \cref{fig:bx-gap} shows the gap of every run, and \cref{tab:bx-stats} gives
the tests.

\input{\bxdir tables/tab_main}

\begin{figure}[tbp]
\centering
\includegraphics[width=\linewidth]{\bxdir figures/fig_gap_box}
\caption{Gap of each run to the reference value (symmetric log scale; linear between $-1$ and $1$\,\%).
Boxes show quartiles and whiskers the range over the feasible runs; every run is also drawn as a point.
Construct and MINLP are single runs; ``0/1'': no feasible layout. Annotation above a box: number of feasible
runs when it is below 10. References: O7, the proven optimum (MINLPLib~\cite{minlplib_o7}); vC10a, SC30 from
Jeong and Chae~\cite{jeong2026}; AB20 from Gon\c{c}alves and Resende~\cite{goncalves2015}; Du62 from Kang and
Chae~\cite{kang2017} as cited in~\cite{jeong2026}. The polish rows are in \cref{tab:bx-main}.}
\label{fig:bx-gap}
\end{figure}

\input{\bxdir tables/tab_stats}

\paragraph{Feasibility.} All heuristics found feasible layouts in every run, with one exception: Swap found a
feasible layout on Du62 in only \bxNfeasSwapDusixtytwo{} of \bxNrunsSwapDusixtytwo{} runs. The MINLP found the
optimum of O7 (\bxMinlpOsevenCostTen, first reached after \bxMinlpOsevenTfirstOpt~s) but could not prove it within
300~s (gap \bxMinlpGapOseven\,\%). On vC10a and SC30 its best layouts are far from the others (\bxMinlpIncumbentvCtena{}
and \bxMinlpIncumbentSCthirty), and on AB20 and Du62 it found no feasible layout. Its lower bounds after 300~s are
weak (for example \bxMinlpBoundABtwenty{} on AB20).

\paragraph{Solution quality.} Our method has the lowest mean cost on all five instances; on vC10a, SA, CRO and their
polished versions have the same mean (every run of these methods ends at the published layout,
\bxBestOursvCtena). On O7 every run of our method reaches the optimum up to our area safety margin, while Swap, SA
and CRO end at the best slicing layout 131.69 in every run; one polish of their layouts closes most of this gap (mean gap
\bxGapMeanThreeSAPolishOseven\,\% for SA + polish against \bxGapMeanThreeOursOseven\,\% for our method,
\cref{tab:bx-main}). Against the strongest baseline, SA, the difference is significant on SC30 (adjusted
$p=\bxPholmSASCthirty$) but not on AB20 (adjusted $p=\bxPholmSAABtwenty$) or Du62 (adjusted $p=\bxPholmSADusixtytwo$);
the same holds against SA + polish (\cref{tab:bx-stats}). Against Swap, GA, CRO and CRO + polish, the differences are
significant on every instance where the costs differ. The construction heuristic is between
\bxGapMeanConstructOseven\,\% (O7) and \bxGapMeanConstructSCthirty\,\% (SC30) above the references.

\paragraph{The polish control.} One SOCP solve improves the SA layouts on O7 and SC30 in every run (median
reduction \bxPolishMedRedSAPolishSCthirty\,\% on SC30) and changes almost nothing on the full floors (AB20: an
improvement in \bxPolishNimpSAPolishABtwenty{} of 10 runs; Du62: median \bxPolishMedRedSAPolishDusixtytwo\,\%).
On SC30, SA + polish has mean gap \bxGapMeanSAPolishSCthirty\,\%, against \bxGapMeanSASCthirty\,\% for SA and
\bxGapMeanOursSCthirty\,\% for our method. So part of our advantage on SC30 comes from the SOCP geometry itself, and
part from the search.

\paragraph{Anytime behaviour.} \Cref{fig:bx-conv} shows the median best feasible cost over time. Our method has
a feasible layout early and improves steadily. SA has a feasible layout in most runs only in the last part of the
budget, and then improves fast. At the end of the budget its median gap is \bxGapMedianSAABtwenty\,\% on AB20,
\bxGapMedianSASCthirty\,\% on SC30 and \bxGapMedianSADusixtytwo\,\% on Du62, against
\bxGapMedianOursABtwenty\,\%, \bxGapMedianOursSCthirty\,\% and \bxGapMedianOursDusixtytwo\,\% for our method.
Swap and GA are not shown because their runs were made with a code version that did not record the anytime trace.

\begin{figure}[tbp]
\centering
\includegraphics[width=\linewidth]{\bxdir figures/fig_convergence}
\caption{Anytime behaviour: gap of the best feasible cost found so far, median over 10 seeds (line) and
interquartile range (band; drawn once all 10 runs have a feasible layout). The median is drawn once more than
half of the runs have a feasible layout. Dashed line: reference value (AB20 \bxRefABtwenty, SC30
\bxRefSCthirty, Du62 \bxRefDusixtytwo). The marker shows the median at the end of the budget.
Swap and GA have no recorded traces in this data version and are not shown. The MINLP incumbent would be drawn as
a separate step line, but it has no feasible layout on AB20 and Du62, and on SC30 its incumbent stays above the
plotted range. Upper parts of the curves are cut off.}
\label{fig:bx-conv}
\end{figure}

\paragraph{Cost of Stage~2.} The SOCP that evaluates one sequence pair took on average
\bxSocpMsOseven~ms on O7, \bxSocpMsABtwenty~ms on AB20, \bxSocpMsSCthirty~ms on SC30 and
\bxSocpMsDusixtytwo~ms on Du62. On the full floors many changed sequence pairs have no feasible geometry:
\bxSocpInfPctABtwenty\,\% of the SOCPs were infeasible on AB20, against \bxSocpInfPctSCthirty\,\% on SC30,
which has free area. On Du62 the budget rule applied in all \bxPolishOnlyDusixtytwo{} runs, so Stage~2 was a
single polish there. As a consequence, the full method and the two Stage-2 ablations
(\cref{sec:bx-ablation}) gave identical results on Du62 in every seed.

\subsection{Best layouts and published values}
\label{sec:bx-records}

\Cref{fig:bx-layouts} shows the best layout of our method for each instance and \cref{tab:bx-records} compares
these layouts with the best published values. The best layouts come from longer, separately configured runs (up
to \bxRecMaxTL~s, with other settings of the time split and the elite pool); they are not protocol results. The
column ``Best protocol run'' gives the lowest cost over all protocol runs of all methods. All layouts were checked
again by the evaluator.

\input{\bxdir tables/tab_records}

\begin{figure}[tbp]
\centering
\includegraphics[width=\linewidth]{\bxdir figures/fig_layouts}
\caption{Best layouts of our method (record runs, not protocol runs), drawn to scale within each panel (thick
outline: floor). Numbers are department indices; lines join the centroids of the department pairs with the largest
flows (top 15\,\% of the pairs with non-zero flow, at least 5 and at most 40; 25 for Du62), with line width growing
with the flow. White areas in SC30 are free floor. Costs: O7 \bxRecBestOsevenSeven, vC10a \bxRecBestvCtena, AB20
\bxRecBestABtwenty, SC30 \bxRecBestSCthirty, Du62 \bxRecBestDusixtytwo.}
\label{fig:bx-layouts}
\end{figure}

Published values were obtained under different rules, so we compare only where the rules match, and we
report the others as they are.
\begin{itemize}
\item \textbf{O7: no record is possible.} The instance is solved to proven optimality (\bxRefOsevenFull).
The value 131.69 of Jeong and Chae~\cite{jeong2026} is the best slicing layout. The value 131.56 of
Gon\c{c}alves and Resende~\cite{goncalves2015} lies below the proven optimum, so their layout cannot meet every area
exactly; it uses a tangent outer approximation of the area constraint. Our layout costs \bxRecBestOsevenSeven, which
is the optimum up to our area safety margin (relative difference $\bxRecOsevenRelGap$).
\item \textbf{vC10a.} Our best layout is the published layout of~\cite{jeong2026} and has the same cost,
\bxRecBestvCtena. In the protocol runs it is reached by \bxProtBestNtiedvCtena{} methods (a tie, not a record).
\item \textbf{AB20.} Our best layout (\bxRecBestABtwenty) is \bxDeltaABtwentyGR\,\% below the value of
Gon\c{c}alves and Resende~\cite{goncalves2015}, which uses the area outer approximation and a flow-matrix version
that is not recorded. The lower published value 4{,}927.27~\cite{jeong2026} belongs to a layout in which two
departments have aspect ratio 4.008, above $\alpha=4$, so it is not feasible under the strict limit. The plain SA
baseline also goes below that value (best \bxBestSAABtwenty). No published AB20 value that we know is both
exact-area and strictly feasible.
\item \textbf{SC30.} Our best layout (\bxRecBestSCthirty) is \bxDeltaSCthirtyJC\,\% below the best value of
Jeong and Chae~\cite{jeong2026}. The comparison is in our favour: their layouts are slicing layouts that fill the
free area with dummy blocks, while our model allows non-slicing layouts and free space anywhere. Every layout of
theirs is feasible for us, but not the other way round.
\item \textbf{Du62.} Our best layout (\bxRecBestDusixtytwo, one long Stage-1 anneal and one SOCP polish) is
\bxDeltaDusixtytwoJC\,\% below the value of Jeong and Chae~\cite{jeong2026} on the same square floor, and
\bxDeltaDusixtytwoKC\,\% below the value cited for Kang and Chae~\cite{kang2017}. That value is a secondary citation,
and its floor is not confirmed. The plain SA baseline also goes below both values in a protocol run (best
\bxBestSADusixtytwo, \bxDeltaDusixtytwoJCsa\,\% and \bxDeltaDusixtytwoKCsa\,\% below them), so the improvement over
the literature on Du62 comes from slicing-tree annealing, not from Stage~2.
\end{itemize}

\subsection{Ablation}
\label{sec:bx-ablation}

\Cref{tab:bx-ablation} and \cref{fig:bx-ablation} compare the full method with five variants under the same
protocol, and \cref{fig:bx-stage2} shows how much Stage~2 lowers the cost within each run.

\input{\bxdir tables/tab_ablation}

\begin{figure}[tbp]
\centering
\includegraphics[width=\linewidth]{\bxdir figures/fig_ablation}
\caption{Ablation on AB20, SC30 and Du62: mean gap ($\pm$ one standard deviation, 10 seeds) and best gap of the
full method and five variants. References as in \cref{fig:bx-gap}. Note the different horizontal scales.}
\label{fig:bx-ablation}
\end{figure}

\begin{figure}[tbp]
\centering
\includegraphics[width=0.6\linewidth]{\bxdir figures/fig_stage2}
\caption{Cost reduction by Stage~2 within each run, $100\,(c_1-c)/c_1$, where $c_1$ is the best Stage-1 cost
of the run and $c$ the final cost. Triangles: Stage~1 + polish (one SOCP solve on the final Stage-1 layout);
circles: full method (Stage~1 for half of the budget, then Stage~2). Bars: medians. Symmetric log scale, linear
below $10^{-3}$\,\%. On Du62 the full method also reduces to a polish (budget rule).}
\label{fig:bx-stage2}
\end{figure}

\paragraph{Stage~1 is needed.} Without Stage~1, Stage~2 starts from a random feasible sequence pair. The mean gap
rises to \bxGapMeanOursNoSoneABtwenty\,\% on AB20 and \bxGapMeanOursNoSoneSCthirty\,\% on SC30, and on Du62 no
random feasible sequence pair was found in any run (\bxNfeasOursNoSoneDusixtytwo/\bxNrunsOursNoSoneDusixtytwo).

\paragraph{Stage~2 matters on SC30 and O7 only.} On SC30 a single polish lowers the cost by a median of
\bxStageTwoMedRedOursStwoSCthirty\,\% and the full Stage~2 by \bxStageTwoMedRedOursSCthirty\,\%
(\cref{fig:bx-stage2}); the mean gap falls from \bxGapMeanOursSoneSCthirty\,\% (Stage~1 only) to
\bxGapMeanOursStwoSCthirty\,\% (polish) and \bxGapMeanOursSCthirty\,\% (full). The full method is significantly better
than both variants on SC30 (adjusted $p=\bxPholmOursSoneSCthirty$ and $\bxPholmOursStwoSCthirty$,
\cref{tab:bx-stats}). On O7 a single polish already turns the best slicing layout into the optimal non-slicing
layout. On the full floors the effect is small or absent. On vC10a and AB20 the polish lowers the cost by less than
$10^{-5}$\,\% in every run, because a fixed sequence pair on a full floor fixes the geometry. On AB20 the Stage-2
search lowers the cost by a median of \bxStageTwoMedRedOursABtwenty\,\% relative to the Stage-1 result at half time,
but the full method (mean \bxMeanOursABtwenty) is not better than Stage~1 alone run for the whole budget (mean
\bxMeanOursSoneABtwenty; adjusted $p=\bxPholmOursSoneABtwenty$). On Du62 Stage~2 is a single polish.

\paragraph{Elite pool and neighbour moves.} With an elite pool of size one, the mean is
\bxMeanOursNoEliteABtwenty{} on AB20 and \bxMeanOursNoEliteSCthirty{} on SC30, against \bxMeanOursABtwenty{}
and \bxMeanOursSCthirty{} for the full method; with uniform partner choice it is \bxMeanOursUnifABtwenty{} and
\bxMeanOursUnifSCthirty. These differences are within one standard deviation, so ten seeds do not show an effect
of either component. The best AB20 layout of all protocol runs (\bxBestOursNoEliteABtwenty) was found by the
variant with elite pool one.

\subsection{Threats to validity}
\label{sec:bx-threats}

\begin{itemize}
\item \textbf{Tuning on test instances.} The Stage-1 parameters of our method were tuned on AB20 and SC30, and the
Stage-2 temperatures on SC30; these are benchmark instances. The baselines were not tuned. The best layouts in
\cref{tab:bx-records} come from longer runs with other settings and are not protocol results.
\item \textbf{Wall-clock budgets.} Budgets are in seconds on one machine. Results depend on the machine and on
the load from the other two parallel runs, and runs are not bit-reproducible. The runs of our method overran the
budget by up to \bxOverOursMaxPctDusixtytwo\,\% on Du62 (\cref{sec:bx-protocol}); the baselines did not. The exact
models ran on a different machine with two cores.
\item \textbf{Baselines.} Swap, GA and SA are our own implementations on the STS encoding, with parameters
chosen by us; GA has no local search. CRO is our reproduction, run under the common budget instead of the
paper's stopping rule and with a static penalty on SC30; it does not reproduce the published SC30 results even
with the full stopping rule. Stronger baselines may exist.
\item \textbf{Exact model.} The MINLP got one 300~s run per instance; with more time or a stronger formulation it
may do better, as the 3600~s O7 run shows.
\item \textbf{Reference values.} Several published values use other area rules, other data versions or an
unconfirmed floor (\cref{tab:bx-records}); gaps to them are indicative only. The reference values were checked
against our literature notes, not against the original papers in this study.
\item \textbf{Sample size.} Ten seeds per method and instance give limited power; the AB20 and Du62 differences
to SA and most ablation differences are not significant.
\item \textbf{Tolerance.} The evaluator accepts violations up to $10^{-6}$. Solver-based layouts (SOCP, MINLP) meet
the limits to about $10^{-7}$ (the best AB20 layout has a largest residual of order $10^{-10}$, since a full floor
leaves no room for a safety margin); slicing layouts are exact to about $10^{-9}$.
\end{itemize}
